\chapter{Coordonnées angulaires conjuguées et décomposition spectrale}
\label{chap:canales}

Le passage de la clôture arithmétique de \(\alpha\) à la hiérarchie
dimensionnelle s'effectue au moyen de deux coordonnées angulaires. La première provient
directement de \(\alphaHMT\) ; la seconde est la phase orientée produite par
la réalisation analytique régionale de la microfibre quintuple. Toutes deux arrivent en
mécanique dimensionnelle déjà construites, avant l'introduction de la carte exponentielle
qui engendrera les moments globaux.

\section{Transfert des coordonnées angulaires dérivées dans HMT}

Les Parties I--III ont construit des dépendances distinctes ayant une racine commune dans
\(\alpha_{\rm HMT}\). Le caractère pré-retour provient de
\[
(\alpha_{\rm HMT},\varphi;\,9,5,4,12,54,\text{ règle graduée})
\longrightarrow H_5,
\]
tandis que la coordonnée angulaire et le déterminant proviennent de
\[
\alpha_{\rm HMT}\longrightarrow
A=1000\alpha_{\rm HMT}\,{}^\circ,
\qquad
(\alpha_{\rm HMT},\pi)\longrightarrow
D_A=e^{-100\pi\alpha_{\rm HMT}/9}.
\]
La microfibre régionale fixe, par une dépendance distincte,
\(\mathcal F_\pi\to\rho_\pi=169\), et avec elle
\[
\mathcal C_\pi(\alpha_{\rm HMT})
=169\alpha_{\rm HMT}^{6}
+\frac{D_A\alpha_{\rm HMT}^{7}}
       {1-D_A\alpha_{\rm HMT}}.
\]
La branche de \(H_5\) et la correction régionale ne confluent qu'à l'étape
suivante :
\[
\begin{aligned}
y_*&=\frac{\pi}{90}(H_5+\alpha_{\rm HMT})
     -\mathcal C_\pi(\alpha_{\rm HMT}),\\
C^*&=\frac{180}{\pi}y_*.
\end{aligned}
\]
Désormais, \(A\) désigne la coordonnée angulaire directe et \(C^*\) la
coordonnée angulaire régionale conjuguée ; ces dénominations fixent leurs types
sans introduire une seconde règle de sélection.
Par conséquent, \(H_5\) n'est pas une sortie de \(A\). Ici, \(y_*\) est la coordonnée
analytique en radians et \(C^*\) sa coordonnée sexagésimale. La mécanique
dimensionnelle reçoit ces sorties déjà déterminées ; elle ne les redéfinit pas.
La coordonnée \(A\) n'est pas une entrée directe de la formule de \(C^*\) : les deux
objets sont couplés pour la première fois dans \((A,C^*)\mapsto q_\pm\). L'égalité
\(\alpha_{\rm HMT}=A/1000\) n'est que le changement de coordonnée inverse.
Dans cette Partie, nous conservons les notations
\[
\begin{aligned}
 A&=7.297352569283800997285105472380663^\circ,\\
 \Cstar&=2.123738338968645724261951380393\ldots{}^\circ.
\end{aligned}
\]
L'arrondi utilisé dans les tableaux est
\(2.123738338968646^\circ\). L'entier \(169\), le déterminant \(D_A\) et
la récurrence régionale appartiennent à la démonstration précédente ; aucune
comparaison SI ni la fonctionnelle hexade--octade n'interviennent dans leur génération.

L'identité
angulaire
\[
 \etaRet=\frac{\Cstar}{2}-\alphaAngle,
 \qquad
 \alphaAngle:=\alphaHMT\,{\rm deg},
\]
définit la coordonnée angulaire adimensionnelle postérieure au retour ; elle n'est pas utilisée
pour engendrer \(\Cstar\). De manière équivalente,
une fois la construction régionale effectuée,
\[
 \Cstar=2(\etaRet+\alphaAngle).
\]
Cette dernière égalité est la forme inverse de la même relation et non une
seconde règle de sélection de \(C^*\).

La réalisation régionale corrige \(H_5\) avant de produire \(\etaRet\) ; par
conséquent, \(H_5\),
\(\etaRet\) et \(\Cstar\) ont des types mathématiques distincts. Le symbole
\(\hbarHMT\) est réservé à la grandeur sur la droite d'action obtenue
seulement après application de l'échelle \(10^{-34}\mathcal U_S\). La comparaison
avec une mantisse du SI appartient au contrôle métrologique externe et laisse
invariante la construction précédente.

\section{Phases internes et carte exponentielle}

Les phases abstraites associées aux deux angles sont
\[
 a=\frac A{360},\qquad c=\frac{\Cstar}{360}
 \quad\text{dans }\RR/\ZZ.
\]
Avant d'appliquer une exponentielle réelle, on choisit les représentants canoniques
\(\widetilde a,\widetilde c\in[0,1)\). L'exponentielle agit sur ces
représentants, non directement sur les classes du quotient.
La conversion en radians n'est introduite que lors du choix de la carte
archimédienne
\[
 x:=A_{\rm rad}=\frac{\pi A}{180},\qquad
 y:=\Cstarrad=\frac{\pi\Cstar}{180},
\]
\[
 q_-=e^{-x+y},\qquad q_+=e^{-x-y}.
\]
Le facteur \(\pi/180\) appartient à ce changement analytique de coordonnées. La
clôture arithmétique de \(\alpha\), formulée avant cette conversion, est
indépendante de ce choix de carte angulaire.

\section{Algèbre involutive centrale}

Considérons la représentation régulière bidimensionnelle de l'algèbre réelle
engendrée par une involution. Dans une base spectrale, nous fixons
\[
 R=\operatorname{diag}(1,-1),\qquad I=I_2,
\]
de sorte que les deux sous-espaces propres ont une multiplicité un. Soient
\[
 P_+=\frac{I+R}{2},\qquad P_-=\frac{I-R}{2}
\]
ses projecteurs spectraux. L'élément central déterminé par les
coordonnées angulaires est
\[
 B_{A,\Cstar}=AI+\Cstar R
 =(A+\Cstar)P_++(A-\Cstar)P_-.
\]
L'involution \(\Cstar\mapsto-\Cstar\) conserve les projecteurs fixes
\(P_\pm\) et échange les deux coefficients spectraux ---de manière équivalente,
les canaux associés---. En appliquant l'exponentielle, on obtient
\[
 T=e^{-xI-yR}=q_+P_++q_-P_-.
\]
Par conséquent, les deux canaux sont les valeurs propres d'un unique élément de
l'algèbre engendrée par \(I\) et \(R\).

\section{Inversibilité de la carte des canaux}

\begin{theorem}[Reconstruction des coordonnées angulaires]
Para \(q_\pm>0\),
\[
 x=-\frac12\log(q_+q_-),
 \qquad
 y=\frac12\log\frac{q_-}{q_+}.
\]
Par conséquent, l'application
\((A,\Cstar)\mapsto(q_+,q_-)\) est injective. L'échange
\(q_+\leftrightarrow q_-\) conserve \(A\) et transforme
\(\Cstar\mapsto-\Cstar\).
\end{theorem}
\begin{proof}
La multiplication et le quotient des deux exponentielles donnent
\(q_+q_-=e^{-2x}\) et \(q_-/q_+=e^{2y}\). En prenant les logarithmes, on retrouve
\(x\) et \(y\). Comme \(\pi/180\neq0\), les deux angles sont déterminés.
\end{proof}

\section{Identités d'addition et de récurrence}

Dans cette représentation bidimensionnelle, les moments pair et impair de l'opérateur
\(T\) sont
\[
 c_n=\Tr(T^n)=q_-^n+q_+^n,
 \qquad
 s_n=-\Tr(RT^n)=q_-^n-q_+^n.
\]

L'hypothèse de multiplicité un fait partie du domaine de ces formules.
Dans une représentation de multiplicités \(m_+\) et \(m_-\), les traces
correspondantes seraient
\(m_+q_+^n+m_-q_-^n\) et
\(m_-q_-^n-m_+q_+^n\), respectivement.

\begin{theorem}[Lois de composition des moments]
Pour \(m,n\geq0\),
\[
 c_{m+n}=\frac{c_mc_n+s_ms_n}{2},
 \qquad
 s_{m+n}=\frac{s_mc_n+c_ms_n}{2},
\]
\[
 c_n^2-s_n^2=4e^{-2nx},
\]
et les deux suites satisfont
\[
 X_{n+2}=c_1X_{n+1}-e^{-2x}X_n.
\]
De plus,
\[
 \partial_Ac_n=-\frac{n\pi}{180}c_n,
 \quad
 \partial_As_n=-\frac{n\pi}{180}s_n,
\]
\[
 \partial_{\Cstar}c_n=\frac{n\pi}{180}s_n,
 \quad
 \partial_{\Cstar}s_n=\frac{n\pi}{180}c_n.
\]
\end{theorem}
\begin{proof}
On développe les produits de
\(q_-^m\pm q_+^m\) avec \(q_-^n\pm q_+^n\). L'identité quadratique résulte
de \((u+v)^2-(u-v)^2=4uv\). La récurrence correspond au polynôme
caractéristique \(z^2-c_1z+q_-q_+\), et les dérivées s'obtiennent par la
règle de la chaîne.
\end{proof}

\begin{corollary}[Détermination de la hiérarchie complète]
Une fois \(A\) et \(\Cstar\) fixés, tous les couples \((c_n,s_n)\) sont
déterminés par la récurrence. Aucune espèce physique n'introduit de canal
supplémentaire.
\end{corollary}

Enfin,
\[
 s_n=2e^{-nx}\sinh(ny),
 \qquad
 c_n=2e^{-nx}\cosh(ny).
\]
La coordonnée \(A\) détermine la décroissance commune et \(\Cstar\) la séparation
orientée. L'ordre décimal de la droite d'action n'est pas introduit dans ces
formules. Il se déduit dans le \cref{chap:escala-accion} à partir du conoyau
entier du bloc local de Hadamard. Le choix ultérieur d'une base
exprimée en \(\mathrm{J\,s}\) appartient au transducteur métrologique et ne
modifie pas les identités adimensionnelles.
